\section{Literature Review}
Find roughly 10–15 strong studies on bullying and educational achievement.
Prioritise PIRLS/TIMSS studies, South African evidence, and comparable middle-income countries.
Record each paper’s bullying measure, outcome, controls, method, and main result.
Use this to construct a simple conceptual framework and select covariates

•	bullying and academic achievement;
•	bullying and reading achievement specifically;
•	PIRLS or TIMSS bullying studies;
•	South African bullying research;
•	Brazilian or comparable middle-income-country evidence;
•	gender or socioeconomic differences in bullying effects.
Create a literature matrix with:
Study	Country/data	Bullying measure	Outcome	Controls	Method	Main result	Limitation
The literature review should answer:
1.	Why might bullying affect reading achievement?
2.	Which variables might confound the relationship?
3.	Is there evidence that the relationship differs by gender or SES?
4.	Which variables might be consequences of bullying rather than confounders?
That last question is important. Variables such as absenteeism, school belonging, fear, or engagement may be mechanisms through which bullying affects achievement. Automatically controlling for them could remove part of the relationship you are trying to estimate.


França ()